\originalchapter{8.04 作业10：径向谱与氢原子}
以下按MIT原件中文转述, 解答均为\textbf{AI独立解答}; 保留原件的改编信用, 依完整MIT任务独立推导. 来源为Spring 2016 Assignment 10, 2016年4月29日发布, 5月6日12时截止.
\originalsection{第1题：虚波数与束缚极点}
\begin{exercise}有限程势的外区散射解$\psi_k=\e^{\ii\delta}\sin(kx+\delta)$, $x>R$. 解释束缚态为何对应$A_s=\e^{\ii\delta}\sin\delta$在$k=\ii\kappa,\kappa>0$的极点.\end{exercise}
\begin{solution}令$S=\e^{2\ii\delta}=1+2\ii A_s$, 则$\psi_k=(S\e^{\ii kx}-\e^{-\ii kx})/(2\ii)$. 延拓至$k=\ii\kappa$后第一项衰减, 第二项增长. 常规入射归一解的增长项系数固定, 不能直接成为正规束缚态; 在$S$的极点附近将整态乘$(k-\ii\kappa)$再取留数, 增长项消失, 衰减项留下非零解. 反向, 从满足原点正规边界的解$\chi_k=A(k)\e^{\ii kx}+B(k)\e^{-\ii kx}$出发, 束缚时$B(\ii\kappa)=0$, $A(\ii\kappa)\ne0$（否则解处处0）, 而$S=-A/B$. 对孤立简单零点便是极点. 这就是正规束缚态的极点解释, 不是将实轴上的实相移原封不动代成虚数. 适用于解析延拓可成立的短程势、孤立束缚态; 阈值态需另处理.\end{solution}
\originalsection{第2题：非简并共同本征态}
\begin{exercise}厄米$A,B$对易且$A$谱不简并, 证明$A$每个本征态也是$B$本征态.\end{exercise}
\begin{solution}若$Au=au$, 则$A(Bu)=B(Au)=aBu$. 一维本征空间使$Bu=bu$, 也包括$Bu=0$的$b=0$情形. $B$厄米保证$b$实. 有简并时仍保持本征子空间, 但须在其内部再对角化$B$.\end{solution}
\originalsection{第3题：球谐展开与角动量测量}
\begin{exercise}
原件署Based on Ohanian Ch.7 problem17. 径向函数满足$\int_0^\infty r^2|f(r)|^2\dd r=1$, 波函数$\psi=\tfrac14\sqrt{5/\pi}\sin^2\theta(1+\sqrt{14}\cos\theta)\cos2\phi\,f(r)$. (a)球谐展开, 求$L^2,L_z$各读数与概率. (b,c)求均值与不确定度. $\cos2\phi$在整个括号外, 以实际PDF为准.
\end{exercise}
\begin{solution}
使用$Y_{2,\pm2}=\sqrt{15/(32\pi)}\sin^2\theta\e^{\pm2\ii\phi}$和$Y_{3,\pm2}=\sqrt{105/(32\pi)}\sin^2\theta\cos\theta\e^{\pm2\ii\phi}$, 得$\psi/f=(Y_{22}+Y_{2,-2})/\sqrt6+(Y_{32}+Y_{3,-2})/\sqrt3$. 四项概率和$2/6+2/3=1$. $L^2$读数$6\hbar^2,12\hbar^2$概率$1/3,2/3$, $L_z=\pm2\hbar$各$1/2$. 因而$\langle L^2\rangle=10\hbar^2$, $\langle(L^2)^2\rangle=108\hbar^4$, $\Delta(L^2)=2\sqrt2\hbar^2$. $\langle L_z\rangle=0$, $\langle L_z^2\rangle=4\hbar^2$, $\Delta L_z=2\hbar$. $L^2$的不确定度不能误写成$L$的不确定度.
\end{solution}
\originalsection{第4题：球形井}
\begin{exercise}(a)$r<a$内$V=0$、外无限的球井, 求$\ell=0$径向$u$与能谱并对照一维井. (b)改井内$-V_0$、外0, 证$V_0a^2<\pi^2\hbar^2/(8m)$时无束缚态.\end{exercise}
\begin{solution}写$\psi=u(r)Y_{00}/r$, 原点正规要求$u(0)=0$. 无限井$u=\sqrt{2/a}\sin(n\pi r/a)$, $E_n=n^2\pi^2\hbar^2/(2ma^2)$, 恰对应长度$a$一维无限井, 不是长度$2a$. 有限井的$s$波内$\sin kr$, 外$\e^{-\kappa r}$, 接续为$ka\cot ka=-\kappa a$及$(ka)^2+(\kappa a)^2=2mV_0a^2/\hbar^2$. 第一根在$ka>\pi/2$, 阈值为$V_0a^2=\pi^2\hbar^2/(8m)$; 等号对应不归一零能阈值, 实束缚需严格大于. 更高$\ell$多正的离心项, 不会先产生比$s$波更低的束缚态.\end{solution}
\originalsection{第5题：运动氢原子的联合概率与相位}
\begin{exercise}
原件署Based on Ohanian. 质心$\boldsymbol X$, 内坐标$\boldsymbol x=\boldsymbol x_e-\boldsymbol x_p$. 总动量$\pm\boldsymbol p_0$各半, 内态$\phi_{100},\phi_{211}$概率$1/4,3/4$, 两种测量概率不相关. (a)写允许任意相位且略去整体相位的纯态; (b)求总能量均值.
\end{exercise}
\begin{solution}
用盒正规质心动量态$\chi_\pm=L^{-3/2}\e^{\pm\ii\boldsymbol p_0\cdot\boldsymbol X/\hbar}$（两动量不同）, 最一般相符纯态为
\[
\psi=\sum_{s=\pm}\chi_s(\boldsymbol X)
\left[\frac{\e^{\ii\theta_{s1}}}{\sqrt8}\phi_{100}(\boldsymbol x)
+\sqrt{\frac38}\e^{\ii\theta_{s2}}\phi_{211}(\boldsymbol x)\right].
\]
四相位扣除一个整体相位后仍有三个自由相位. 概率不相关只约束四个系数的模; 若相位差依赖$s$, 仍可能纠缠, 不应只写一个乘积态. 令$M=m_e+m_p$, $\mu=m_em_p/M$, $E_n=-R_\mu/n^2$, $R_\mu=\mu e^4/(2\hbar^2)$, 四项在总能量中正交, 得$\langle H\rangle=p_0^2/(2M)-R_\mu(1/4+3/16)=p_0^2/(2M)-7R_\mu/16$. 若$p_0=0$两标签合并, 不再是两个不同动量态, 需另按一个质心态归一.
\end{solution}
\originalsection{第6题：维里定理的氢应用}
\begin{exercise}定态为何使$\mathrm d\langle\Omega\rangle/\mathrm dt=0$? 用$\Omega=\boldsymbol r\cdot\boldsymbol p$证氢$\langle T\rangle=-\langle V\rangle/2$, 求$n$态的$rms$速率与$c$之比, 推广核电荷$Z$并求$\langle1/r\rangle$.\end{exercise}
\begin{solution}定态只乘整体相位, 固定算符的均值不变. 具体$[\Omega,T]=2\ii\hbar T$, $[\Omega,V]=-\ii\hbar\boldsymbol r\cdot\nabla V$, 所以$\dot\Omega=[\Omega,H]/(\ii\hbar)=2T-rV^\prime(r)$; Coulomb势$V=-e^2/r$满足$rV'=-V$, 因而$2\langle T\rangle=-\langle V\rangle$, $\langle T\rangle=-E_n$. $E_n=-me^4/(2\hbar^2n^2)$时$\sqrt{\langle v^2\rangle}/c=\alpha/n$, 核电荷$Z$则为$Z\alpha/n$, 基态$Z\alpha$; 当它接近1非相对论模型失效. 从$\langle V\rangle=2E_n$得$\langle1/r\rangle=1/(a_0n^2)$, 电荷$Z$则$Z/(a_0n^2)$. 本题用Gaussian电磁单位的$e^2$, SI应替成$e^2/(4\pi\epsilon_0)$.\end{solution}
\originalsection{第7题：氢基态、核内概率与正电子素}
\begin{exercise}求氢基态$\langle r\rangle,\langle r^2\rangle$及最可能半径; 核半径1fm时电子核内概率; 正电子素能级与尺寸.\end{exercise}
\begin{solution}径向概率密度$P_r=4r^2\e^{-2r/a_0}/a_0^3$, 用$\int_0^\infty r^s\e^{-br}\dd r=s!/b^{s+1}$得$\langle r\rangle=3a_0/2$, $\langle r^2\rangle=3a_0^2$; 求导$P_r$得峰$r=a_0$. 这是球壳概率峰, 位置体密度则在原点最大. $R=1$fm远小于$a_0=5.29177\times10^{-11}$m, 核内概率近似$4R^3/(3a_0^3)\simeq9.00\times10^{-15}$, 相对误差$O(R/a_0)$. 正电子素约化质量$m/2$, 故$E_n=-6.8\,\mathrm{eV}/n^2$, Bohr尺度$2a_0$, 暂不含湮灭和精细结构.\end{solution}
